% !TeX root = part4.tex

\documentclass[../final.tex]{subfiles}

\begin{document}

% При отдельной сборке продолжаем нумерацию полного конспекта.
\ifSubfilesClassLoaded{\setcounter{chapter}{3}\setcounter{section}{6}}{}

% ============================================================
% Лекция 22.09
% Аппроксимация и интерполяция
% ============================================================

\chapter{Аппроксимация и интерполяция}

\rsubtitle{Интерполяционные полиномы в форме Лагранжа и Ньютона}

\captionsetup{font=footnotesize,labelfont=bf,skip=5pt}

% ============================================================
\section{Задача интерполяции}
% ============================================================

Пусть функция $y(x)$ задана в $n+1$ точках, то есть известны значения
%
\begin{equation*}
    y(x_0)=y_0,\qquad y(x_1)=y_1,\qquad\ldots,\qquad y(x_n)=y_n.
\end{equation*}
%
Точки $x_0,x_1,\ldots,x_n$ называются \rconcept{узлами интерполяции}.
Требуется построить непрерывную интерполирующую функцию $f(x)$,
совпадающую с заданной функцией во всех узлах:
%
\begin{equation*}
    \begin{cases}
        f(x_0)=y_0,\\
        f(x_1)=y_1,\\
        \hspace{1.5em}\vdots\\
        f(x_n)=y_n.
    \end{cases}
\end{equation*}

В задаче \rconcept{алгебраической интерполяции} в качестве $f(x)$
ищем полином степени не выше $n$:
%
\begin{equation*}
    P_n(x)=a_0+a_1x+a_2x^2+\cdots+a_nx^n,
\end{equation*}
%
удовлетворяющий условиям
%
\begin{equation}
    P_n(x_i)=y_i,\qquad i=0,1,\ldots,n.
    \label{eq:ovf_interpolation_conditions}
\end{equation}

% ============================================================
\section{Существование и единственность интерполяционного полинома}
% ============================================================

\textbf{Теорема.}
Если все узлы попарно различны, то есть
$x_i\ne x_j$ при $i\ne j$, задача алгебраической интерполяции
имеет единственное решение среди полиномов степени не выше $n$.

\subsection{Система для коэффициентов}

Подставляя полином в условия~\eqref{eq:ovf_interpolation_conditions},
получаем $n+1$ линейное уравнение для $n+1$ неизвестного коэффициента:
%
\begin{equation*}
    \begin{cases}
        a_0+a_1x_0+a_2x_0^2+\cdots+a_nx_0^n=y_0,\\
        a_0+a_1x_1+a_2x_1^2+\cdots+a_nx_1^n=y_1,\\
        \hspace{6em}\vdots\\
        a_0+a_1x_n+a_2x_n^2+\cdots+a_nx_n^n=y_n.
    \end{cases}
\end{equation*}
%
В матричной форме эта система записывается как
%
\begin{equation*}
    \underbrace{
    \begin{pmatrix}
        1&x_0&x_0^2&\cdots&x_0^n\\
        1&x_1&x_1^2&\cdots&x_1^n\\
        \vdots&\vdots&\vdots&\ddots&\vdots\\
        1&x_n&x_n^2&\cdots&x_n^n
    \end{pmatrix}}_{A}
    \begin{pmatrix}a_0\\a_1\\\vdots\\a_n\end{pmatrix}
    =\begin{pmatrix}y_0\\y_1\\\vdots\\y_n\end{pmatrix}.
\end{equation*}
%
Она имеет единственное решение тогда и только тогда, когда $\det A\ne0$.

\subsection{Определитель Вандермонда}

\rconcept{Определитель Вандермонда} равен
%
\begin{equation}
    V_n(x_0,\ldots,x_n)
    =\begin{vmatrix}
        1&x_0&\cdots&x_0^n\\
        1&x_1&\cdots&x_1^n\\
        \vdots&\vdots&\ddots&\vdots\\
        1&x_n&\cdots&x_n^n
    \end{vmatrix}
    =\prod_{0\le j<i\le n}(x_i-x_j).
    \label{eq:ovf_vandermonde}
\end{equation}

Докажем формулу методом математической индукции.
При $n=1$ имеем
%
\begin{equation*}
    V_1(x_0,x_1)=\begin{vmatrix}1&x_0\\1&x_1\end{vmatrix}=x_1-x_0.
\end{equation*}

Предположим, что формула верна для $n=k$, и рассмотрим $n=k+1$.
Последовательно, начиная с последнего столбца, вычтем из каждого
столбца предыдущий, умноженный на $x_0$. Получим
%
\begin{equation*}
    V_{k+1}
    =\begin{vmatrix}
        1&0&0&\cdots&0\\
        1&x_1-x_0&x_1(x_1-x_0)&\cdots&x_1^k(x_1-x_0)\\
        \vdots&\vdots&\vdots&\ddots&\vdots\\
        1&x_{k+1}-x_0&x_{k+1}(x_{k+1}-x_0)&\cdots&x_{k+1}^k(x_{k+1}-x_0)
    \end{vmatrix}.
\end{equation*}

Разложим определитель по первой строке и вынесем из каждой оставшейся
строки множитель $x_i-x_0$:
%
\begin{align*}
    V_{k+1}
    &=\prod_{i=1}^{k+1}(x_i-x_0)
    \begin{vmatrix}
        1&x_1&\cdots&x_1^k\\
        \vdots&\vdots&\ddots&\vdots\\
        1&x_{k+1}&\cdots&x_{k+1}^k
    \end{vmatrix}\\[3pt]
    &=\prod_{i=1}^{k+1}(x_i-x_0)
      \prod_{1\le j<i\le k+1}(x_i-x_j)\\
    &=\prod_{0\le j<i\le k+1}(x_i-x_j).
\end{align*}
%
Формула доказана.

Если все узлы различны, ни один множитель в
\eqref{eq:ovf_vandermonde} не равен нулю. Следовательно,
%
\begin{equation*}
    \det A\ne0,
\end{equation*}
%
а значит, коэффициенты $a_0,\ldots,a_n$ определяются однозначно.
Тем самым доказаны \rresult{существование и единственность} полинома $P_n(x)$.

% ============================================================
\section{Интерполяционный полином в форме Лагранжа}
% ============================================================

Будем искать полином в виде
%
\begin{equation*}
    P_n(x)=\sum_{i=0}^{n}y_i\varphi_i(x),
\end{equation*}
%
где базисные полиномы $\varphi_i(x)$ удовлетворяют условиям
%
\begin{equation*}
    \varphi_i(x_j)=\delta_{ij}
    =\begin{cases}1,&i=j,\\0,&i\ne j.\end{cases}
\end{equation*}

Полином $\varphi_i(x)$ должен обращаться в нуль во всех узлах,
кроме $x_i$. Поэтому положим
%
\begin{equation*}
    \varphi_i(x)=C_i\prod_{\substack{j=0\\j\ne i}}^{n}(x-x_j).
\end{equation*}
%
Из условия $\varphi_i(x_i)=1$ находим
%
\begin{equation*}
    C_i=\frac{1}{\displaystyle\prod_{\substack{j=0\\j\ne i}}^{n}(x_i-x_j)}.
\end{equation*}

Таким образом, \rconcept{интерполяционный полином Лагранжа} имеет вид
%
\begin{equation}
    P_n(x)=\sum_{i=0}^{n}y_i
    \prod_{\substack{j=0\\j\ne i}}^{n}\frac{x-x_j}{x_i-x_j}.
    \label{eq:ovf_lagrange}
\end{equation}
%
Каждый базисный полином имеет степень $n$, поэтому степень их суммы
не превосходит $n$. В узле $x_i$ в сумме остаётся только слагаемое $y_i$.

\remark[Недостаток формы Лагранжа.]
При добавлении нового узла приходится перестраивать все базисные
полиномы. Для последовательного дополнения сетки удобнее форма Ньютона.

% ============================================================
\section{Разделённые разности}
% ============================================================

Пусть известны значения $y(x)$ в попарно различных узлах
$x_0,\ldots,x_n$. Введём \rconcept{разделённые разности}.

Разделённая разность нулевого порядка --- само значение функции:
%
\begin{equation*}
    y^{(0)}(x_i)=y(x_i).
\end{equation*}
%
Разделённая разность первого порядка:
%
\begin{equation*}
    y^{(1)}(x_0,x_1)=\frac{y(x_0)-y(x_1)}{x_0-x_1}.
\end{equation*}
%
Второго порядка:
%
\begin{equation*}
    y^{(2)}(x_0,x_1,x_2)
    =\frac{y^{(1)}(x_0,x_1)-y^{(1)}(x_1,x_2)}{x_0-x_2}.
\end{equation*}
%
В общем случае
%
\begin{equation}
    y^{(n)}(x_0,\ldots,x_n)
    =\frac{y^{(n-1)}(x_0,\ldots,x_{n-1})
          -y^{(n-1)}(x_1,\ldots,x_n)}{x_0-x_n}.
    \label{eq:ovf_divided_difference}
\end{equation}

\remark
Здесь верхний индекс $(n)$ при нескольких аргументах означает порядок
разделённой разности. Производную функции в одной точке ниже будем
записывать как $y^{(n)}(x)$.

\subsection{Явная формула и симметрия}

\textbf{Теорема.} Для попарно различных узлов справедлива формула
%
\begin{equation*}
    y^{(n)}(x_0,\ldots,x_n)
    =\sum_{i=0}^{n}
    \frac{y(x_i)}{\displaystyle\prod_{\substack{j=0\\j\ne i}}^{n}(x_i-x_j)}.
\end{equation*}
%
Следовательно, разделённая разность является
\rresult{симметричной функцией своих аргументов}: её значение
не изменяется при перестановке узлов $x_0,\ldots,x_n$.

% ============================================================
\section{Интерполяционный полином в форме Ньютона}
% ============================================================

\subsection{Последовательное понижение степени}

Существование полинома $P_n(x)$ уже доказано. Разность
$P_n(x)-P_n(x_0)$ обращается в нуль при $x=x_0$, поэтому делится
на $x-x_0$ без остатка. Введём полином степени не выше $n-1$:
%
\begin{equation*}
    P_{n-1}^{(1)}(x,x_0)=\frac{P_n(x)-P_n(x_0)}{x-x_0}.
\end{equation*}
%
Деление полинома можно выполнить по схеме Горнера. Получаем
%
\begin{equation*}
    P_n(x)=P_n(x_0)+(x-x_0)P_{n-1}^{(1)}(x,x_0).
\end{equation*}

На следующем шаге аналогично определим
%
\begin{equation*}
    P_{n-2}^{(2)}(x,x_0,x_1)
    =\frac{P_{n-1}^{(1)}(x,x_0)-P_{n-1}^{(1)}(x_1,x_0)}{x-x_1}.
\end{equation*}
%
Здесь числитель снова делится на знаменатель без остатка, и
%
\begin{equation*}
    P_{n-1}^{(1)}(x,x_0)
    =P_{n-1}^{(1)}(x_1,x_0)+(x-x_1)P_{n-2}^{(2)}(x,x_0,x_1).
\end{equation*}

Продолжая процесс, на $k$-м шаге получаем
%
\begin{equation*}
    P_{n-k}^{(k)}(x,x_0,\ldots,x_{k-1})
    =\frac{
    P_{n-k+1}^{(k-1)}(x,x_0,\ldots,x_{k-2})
    -P_{n-k+1}^{(k-1)}(x_{k-1},x_0,\ldots,x_{k-2})
    }{x-x_{k-1}}.
\end{equation*}
%
После $n$ шагов остаётся постоянный полином $P_0^{(n)}$.
Ещё одна разделённая разность равна нулю:
%
\begin{equation*}
    \frac{P_0^{(n)}(x,x_0,\ldots,x_{n-1})
    -P_0^{(n)}(x_n,x_0,\ldots,x_{n-1})}{x-x_n}=0.
\end{equation*}

\subsection{Формула Ньютона}

Последовательно подставляя полученные выражения, приходим к вложенной
записи полинома:
%
\begin{equation*}
    \begin{aligned}
    P_n(x)=P_n(x_0)+(x-x_0)\Bigl[&P_{n-1}^{(1)}(x_0,x_1)\\
    {}+(x-x_1)\bigl[&P_{n-2}^{(2)}(x_0,x_1,x_2)+\cdots\\
    &{}+(x-x_{n-1})P_0^{(n)}(x_0,\ldots,x_n)\bigr]\Bigr].
    \end{aligned}
\end{equation*}
%
Коэффициенты этой записи --- разделённые разности $P_n$ в узлах.
Поскольку $P_n(x_i)=y(x_i)$, они совпадают с разделёнными разностями $y$.
Получаем \rconcept{интерполяционную формулу Ньютона}:
%
\begin{equation}
    P_n(x)=y(x_0)+\sum_{k=1}^{n}
    y^{(k)}(x_0,\ldots,x_k)\prod_{j=0}^{k-1}(x-x_j).
    \label{eq:ovf_newton_interpolation}
\end{equation}

В развёрнутом виде
%
\begin{align*}
    P_n(x)={}&y(x_0)+y^{(1)}(x_0,x_1)(x-x_0)\\
    &+y^{(2)}(x_0,x_1,x_2)(x-x_0)(x-x_1)+\cdots\\
    &+y^{(n)}(x_0,\ldots,x_n)(x-x_0)\cdots(x-x_{n-1}).
\end{align*}

При добавлении нового узла $x_{n+1}$ достаточно добавить одно слагаемое:
%
\begin{equation*}
    P_{n+1}(x)=P_n(x)+y^{(n+1)}(x_0,\ldots,x_{n+1})
    \prod_{j=0}^{n}(x-x_j).
\end{equation*}
%
Уже вычисленные коэффициенты сохраняются, поэтому форма Ньютона
\rresult{удобна для добавления новых узлов}.

% ============================================================
\section{Погрешность интерполяции}
% ============================================================

Интерполяционное приближение имеет вид $y_n(x)=P_n(x)$.
Разность $y(x)-P_n(x)$ обращается в нуль во всех $n+1$ узлах.
Обозначим
%
\begin{equation*}
    \omega_{n+1}(x)=\prod_{i=0}^{n}(x-x_i).
\end{equation*}
%
Для $x$, не совпадающего с узлами, представим погрешность в виде
%
\begin{equation*}
    y(x)-P_n(x)=\omega_{n+1}(x)r(x).
\end{equation*}
%
Найдём выражение для $r(x)$.

\subsection{Вывод остаточного члена}

Зафиксируем точку $u\notin\{x_0,\ldots,x_n\}$ и введём вспомогательную функцию
%
\begin{equation*}
    g(x,u)=y(x)-P_n(x)-\omega_{n+1}(x)r(u).
\end{equation*}
%
Как функция переменной $x$, она имеет $n+2$ различных нуля:
%
\begin{equation*}
    x_0,x_1,\ldots,x_n,u.
\end{equation*}

Пусть $y\in C^{n+1}(I)$, где отрезок $I$ содержит все эти точки.
Тогда $g$ также непрерывно дифференцируема $n+1$ раз по $x$.
Последовательно применяя теорему Ролля, заключаем, что существует
точка $\xi$ между наименьшей и наибольшей из точек
$x_0,\ldots,x_n,u$, в которой
%
\begin{equation*}
    \frac{\partial^{n+1}g}{\partial x^{n+1}}(\xi,u)=0.
\end{equation*}

Поскольку степень $P_n$ не выше $n$, а старший коэффициент
$\omega_{n+1}$ равен единице,
%
\begin{equation*}
    \frac{\partial^{n+1}g}{\partial x^{n+1}}(x,u)
    =y^{(n+1)}(x)-(n+1)!\,r(u).
\end{equation*}
%
Следовательно,
%
\begin{equation*}
    0=y^{(n+1)}(\xi)-(n+1)!\,r(u),
    \qquad
    r(u)=\frac{y^{(n+1)}(\xi)}{(n+1)!}.
\end{equation*}

Возвращаясь к обозначению $x$ для точки, в которой оценивается ошибка,
получаем \rconcept{остаточный член интерполяции}:
%
\begin{equation}
    y(x)-P_n(x)=\frac{y^{(n+1)}(\xi)}{(n+1)!}
    \prod_{i=0}^{n}(x-x_i).
    \label{eq:ovf_interpolation_remainder}
\end{equation}
%
Точка $\xi$ зависит от $x$. В самих узлах погрешность равна нулю.

\subsection{Оценка погрешности на сетке}

Если на рассматриваемом отрезке
%
\begin{equation*}
    \bigl|y^{(n+1)}(t)\bigr|\le C,
\end{equation*}
%
то из~\eqref{eq:ovf_interpolation_remainder} следует
%
\begin{equation*}
    |y(x)-P_n(x)|\le\frac{C}{(n+1)!}
    \prod_{i=0}^{n}|x-x_i|.
\end{equation*}

Пусть узлы упорядочены и расстояние между соседними узлами имеет
порядок $h$:
%
\begin{equation*}
    x_0<x_1<\cdots<x_n,
    \qquad |x_{i+1}-x_i|\sim h.
\end{equation*}
%
При фиксированном $n$ и $x\in[x_0,x_n]$ каждый множитель $|x-x_i|$
имеет порядок $O(h)$. Если производная ограничена независимо от $h$, то
%
\begin{equation*}
    |y(x)-P_n(x)|=O(h^{n+1}).
\end{equation*}
%
Это оценка при сгущении локальной сетки с фиксированным числом узлов.
Вне интервала $[x_0,x_n]$ произведение расстояний до узлов может быстро
расти, и точность экстраполяции ухудшается.

\noindent\begin{minipage}{\linewidth}
\subsection{Эффект Рунге}

\begin{rbox}[breakable=false,left=18pt,right=18pt,top=14pt,bottom=10pt]{[]}
    \centering
    \begin{tikzpicture}[x=1.0cm,y=0.9cm,
        every node/.style={text=rtext,font=\small}]
        \draw[rtext,line width=0.85pt,-{Stealth[length=2.2mm]}]
            (0,0) -- (12.1,0) node[right] {$x$};
        \draw[rtext,line width=0.85pt,-{Stealth[length=2.2mm]}]
            (0,-1.65) -- (0,2.45) node[left] {$\delta$};
        \draw[rc4,line width=1.35pt]
            (0.15,2.0) .. controls (0.4,1.2) and (0.65,-1.5) .. (1.0,-1.4)
            .. controls (1.5,-1.4) and (1.6,1.55) .. (2.05,1.5)
            .. controls (2.4,1.45) and (2.5,0.0) .. (2.8,0);
        \draw[rc4,line width=1.35pt,domain=2.8:7.6,samples=160]
            plot (\x,{0.045*sin(900*(\x-2.8))});
        \draw[rc4,line width=1.35pt]
            (7.6,0) .. controls (8.6,0.1) and (9.2,0.75) .. (9.65,1.6)
            .. controls (9.9,2.1) and (9.85,-1.4) .. (10.25,-1.35)
            .. controls (10.7,-1.3) and (10.5,1.95) .. (11.0,2.05)
            .. controls (11.55,2.2) and (11.65,-0.8) .. (11.95,-1.1);
        \foreach \xx/\lab in {2.8/{x_0},7.6/{x_n}}{
            \draw[rtext,line width=0.7pt] (\xx,-0.09) -- (\xx,0.09);
            \node[below=5pt] at (\xx,0) {$\lab$};
        }
    \end{tikzpicture}
    \captionof{figure}{Схематическое поведение погрешности
    $\delta(x)=y(x)-P_n(x)$: малые колебания на участке интерполяции
    и рост ошибки за его пределами.}
    \label{fig:ovf_runge_effect}
\end{rbox}

\remark
\rconcept{Эффект Рунге} связан с тем, что при увеличении степени
интерполяционного полинома на равномерной сетке могут усиливаться
колебания и возрастать ошибка, в том числе у концов самого интервала
интерполяции. Оценка $O(h^{n+1})$ при фиксированном $n$ не означает
сходимости при $n\to\infty$ на фиксированном отрезке.
\end{minipage}

\end{document}
